\documentclass[instructions.tex]{subfiles}
\begin{document}

\chapter{Le péché est l'unique obstacle au salut, et aux desseins de Dieu.}


Or, ce que toutes les créatures ne peuvent faire, le péché seul le fait. Je dois donc regarder le péché comme l'unique mal qui soit au monde. Les maladies, les disgrâces, ne sont point de véritables maux: il n'y a que le péché qui est toujours un mal et le plus grand de tous les maux, parce qu'il n'y a que le péché qui nous empêche de posséder Dieu.

Quand je serais dans le ciel, si j'y étais avec le péché, je n'y posséderais pas Dieu; j'aurais donc le malheur des réprouvés, qui sont séparés de Dieu. Et quand je serais en enfer, sans péché, avec la grâce de Dieu, je ne serais pas absolument malheureux, parce que je jouirais du bonheur essentiel des saints, qui possèdent Dieu. Saint Augustin disait donc avec raison qu'il aimerait mieux être dans l'enfer sans péché, que d'entrer dans le ciel avec un seul péché\footnote{Mallem purus à peccato gehennam intrare, quàm peccati sorde pollutus cœlorum regna penetrare. S. Aug.}.

Le péché étant donc le seul obstacle à ma béatitude et à mon salut, combien dois-je le craindre!


\primo Le péché est encore le seul obstacle aux desseins du Créateur. Dieu veut nous sauver, le péché nous damne. Dieu crée les hommes pour sa gloire, et le péché la lui ravit. Dieu voudrait régner dans nos coeurs par son amour, par sa grâce, et le péché y fait régner le démon par l'amour des choses du monde.

Tout ce que Dieu fait au-dehors tend à détruire le péché\footnote{Ut destruatur corpus peccati. Rom.6.}. S'il envoie son Fils sur la terre, c'est pour réparer le péché; s'il nous donne des secours, c'est pour nous armer contre le péché; s'il institue des sacrements, c'est pour nous préserver ou pour nous purifier du péché; s'il établit des pasteurs, c'est pour nous instruire et nous faire éviter le péché; s'il nous couronne dans le ciel, c'est pour avoir vaincu le péché; s'il nous punit, c'est pour l'avoir commis.

Le péché renverse tous les desseins de Dieu. Il anéantit, autant qu'il peut, les mérites et le sang du Rédempteur; rend sans effet les sacrements, les grâces, les instructions, la divine parole, les promesses et les menaces de Dieu.

\secundo Quoi! un Dieu s'applique à détruire le péché dans l'univers, et moi je ne m'applique qu'à faire régner le péché dans mon coeur, à l'établir dans les autres par mes exemples, et à le répandre par mes scandales! Ai-je bien pensé que si le péché est opposé à tous les desseins de sa miséricorde, Dieu le fera servir aux desseins de sa justice, pour la confusion et le châtiment du pécheur?


\section{Résolutions}

Je regarderai désormais le péché comme le seul mal qui soit sur la terre.

\primo Tous les matins je ferai la résolution de n'en point commettre dans la journée.

\secundo Surtout je me précautionnerai contre ceux dans lesquels je tombe le plus facilement.

\prayer{O mon Dieu! inspirez-moi une vive horreur du péché.}

\end{document}
